\documentclass{article}
\usepackage{booktabs, graphicx, xcolor, amsmath, setspace}
\usepackage{pgffor, float}
\usepackage{geometry}
\geometry{a4paper,scale=0.8}
\begin{document}

\section*{Experiment setting}

The M5 dataset contains \(1969\) days. We set the initial expanding-window length to
\(S=52\times 7\times 2=728\) days and the maximum forecast horizon to
\(H=28\) days. This produces \(1241,1240,\ldots,1214\) windows for forecast
horizons \(h=1,2,\ldots,28\), respectively. For each horizon, the final
\(365\) windows are reserved for evaluation.

For each forecast horizon and time series, the out-of-sample residuals from
the first \(365\) windows are used to initialize the residual distribution.
At each subsequent forecast origin, the newly observed residual is added and
the distribution is refitted. After the initialization and evaluation windows
are excluded, \(511,510,\ldots,484\) windows remain for training the QOpt
model at horizons \(h=1,2,\ldots,28\), respectively.

Reconciliation training uses \(1000\) draws for
\(\alpha\in\{0.005,0.025,0.975,0.995\}\) and \(500\) draws for the remaining
quantiles. Validation and test evaluation use \(3000\) and \(5000\) draws,
respectively. Results are reported at
\(\alpha\in\{0.005,0.025,0.165,0.25,0.5,0.75,0.835,0.975,0.995\}\).

\begin{table}
    \centering
    \caption{SPL, Normal distribution}
    \begin{tabular}{llllllllll}
\toprule
 & 0.005 & 0.025 & 0.165 & 0.250 & 0.500 & 0.750 & 0.835 & 0.975 & 0.995 \\
method &  &  &  &  &  &  &  &  &  \\
\midrule
base & \textbf{0.017} & \textcolor{red}{0.063} & \textcolor{red}{0.247} & \textcolor{red}{0.314} & 0.402 & 0.334 & \textcolor{red}{0.265} & \textcolor{red}{0.066} & \textbf{0.017} \\
\midrule
QOpt($\beta=100$) & \textcolor{red}{0.020} & \textbf{0.062} & \textbf{0.244} & \textbf{0.311} & \textcolor{red}{0.392} & \textbf{0.321} & \textbf{0.257} & \textbf{0.066} & \textcolor{red}{0.018} \\
ols & 0.029 & 0.109 & 0.363 & 0.398 & 0.430 & 0.533 & 0.461 & 0.127 & 0.033 \\
wls & 0.039 & 0.090 & 0.268 & 0.329 & 0.400 & \textcolor{red}{0.330} & 0.268 & 0.089 & 0.038 \\
shr & 0.024 & 0.092 & 0.322 & 0.370 & \textbf{0.375} & 0.379 & 0.335 & 0.096 & 0.025 \\
sam & 0.195 & 0.743 & 2.426 & 2.536 & 0.422 & 2.596 & 2.466 & 0.748 & 0.196 \\
\bottomrule
\end{tabular}
\end{table}

\begin{table}
    \centering
    \caption{SPL, SkewStudentT distribution}
    \begin{tabular}{llllllllll}
\toprule
 & 0.005 & 0.025 & 0.165 & 0.250 & 0.500 & 0.750 & 0.835 & 0.975 & 0.995 \\
method &  &  &  &  &  &  &  &  &  \\
\midrule
base & \textbf{0.017} & \textcolor{red}{0.063} & \textcolor{red}{0.252} & \textcolor{red}{0.320} & 0.399 & \textbf{0.324} & \textbf{0.258} & \textbf{0.068} & \textbf{0.018} \\
\midrule
QOpt($\beta=100$) & \textcolor{red}{0.019} & \textbf{0.061} & \textbf{0.246} & \textbf{0.314} & 0.404 & \textcolor{red}{0.326} & \textcolor{red}{0.262} & \textcolor{red}{0.071} & \textcolor{red}{0.020} \\
ols & 0.029 & 0.109 & 0.361 & 0.398 & 0.424 & 0.524 & 0.455 & 0.128 & 0.034 \\
wls & 0.039 & 0.093 & 0.272 & 0.331 & \textcolor{red}{0.398} & 0.328 & 0.267 & 0.087 & 0.035 \\
shr & 0.025 & 0.092 & 0.321 & 0.369 & \textbf{0.374} & 0.376 & 0.332 & 0.097 & 0.026 \\
sam & 0.198 & 0.744 & 2.409 & 2.515 & 0.432 & 2.568 & 2.444 & 0.749 & 0.199 \\
\bottomrule
\end{tabular}
\end{table}

\foreach \dist in {normal,skew}{
\foreach \i in {0,...,27}{
    \pgfmathtruncatemacro{\j}{\i+1}
    \begin{table}[H]
    \centering
    \caption{h=\j, \dist}
    \input{M5_ets_\dist_h\i_outsample.tex}
    \end{table}
}}




\end{document}
